% Preliminary measurements for RQ3. Kept out of main.tex: the manuscript
% currently reports no performance numbers in the body.
\subsection{Preliminary layout measurements}
\label{sec:layout-preliminary}

We ran a preliminary version of RQ3 on one cloud VM
(\code{ecs.u1-c1m2.2xlarge}, Ubuntu~26.04). RustFS used EC~2+2 over four
directories of one 100\,GiB ESSD, so the disks are not independent failure
domains. Each layer was compared with its own baseline in paired ABBA
cycles. Figure~\ref{fig:layout-preliminary} summarizes the results.

\paragraph{Backend.} The RustFS extension admits a requested $B$ and
persists $B$ and $q$ with each object. The measured revision~I reads the
$q$-aligned shard windows of a request in one local batch and tells the
kernel exactly which bytes it will read; revision~M is the preceding
integration of the same extension without the batch. With direct S3
requests over eight cycles, cold Range reads ran $4.2$--$9.3\times$ faster
for $B=4\MiB$ and $24$--$26\times$ faster for $B=64\KiB$, cold whole-object
GETs $11$--$14\%$ faster, and cold requests used half the CPU time. Warm
cells ranged from $0.98$ to $1.51\times$. Two warm cells missed a $5\%$
bound; the larger miss came from one late cycle in which several
unrelated cells dropped together.

\paragraph{Client.} The \sys \code{range-heavy} profile omits the hint
for objects within the inline-admission bound, requests $B=64\KiB$ up to
$L_o=8\MiB$ and $256\KiB$ above, and switches ordinary reads to sparse
fetching of 64\,KiB pages. Sparse fetching alone made sequential scans
issue one Range GET per kernel readahead window, so whole-block background
prefetch is restored for three access shapes: a read that continues the
previous one, a read of at least 128\,KiB, and a slice with at least ten
page misses per second. Through FUSE, against the same binary with hints
and the read policy unset, cold 4\,KiB random reads over a 3.75\,GiB set
ran $23\times$ faster with 8 jobs and $3.6\times$ with one, and 64\,KiB
reads $17\times$. Sequential and fresh-file reads, warm reads, a hot set,
and mixed read/write stayed within $3\%$. A 256\,MiB sequential write took
a fixed extra ${\approx}0.25$\,s ($0.91\times$; $0.97\times$ at 512\,MiB).
With the read policy held fixed, $B=256\KiB$ read freshly written files
exactly as fast as $B=64\KiB$ and delivered only $63$--$78\%$ of its sparse
throughput, so the smaller block itself contributes to the gain.

\paragraph{Limits.} Every comparison verified the persisted layout of its
datasets, and every written file passed a checksummed read-back.
Later in two runs, cold reads of the long-lived datasets slowed to
${\approx}0.6$\,s per request for all configurations, while writes and
reads of new files did not; the cold client cells therefore rest on two
cycles. The two layers were not measured end to end together, and the
setup covers a single node, no compression, and SQLite metadata.

% Preliminary measurements for RQ3 (client--RustFS layout co-design).
% Two panels of log2-scaled ratio bars; ratios above 1 favor the candidate.
\begin{figure*}[t]
\centering
\definecolor{gainc}{HTML}{2A78D6}
\definecolor{lossc}{HTML}{E34948}
\begin{tikzpicture}[x=1cm,y=1cm,font=\footnotesize]
  % Ratio r is drawn at x = s * log2(r / 0.5).
  \def\s{0.62}
  \def\barh{0.32}
  \pgfmathsetmacro{\xone}{\s*1}
  \tikzset{
    gain/.style={fill=gainc,draw=none},
    loss/.style={fill=lossc,draw=none},
    flat/.style={fill=black!25,draw=none},
    lbl/.style={anchor=east,inner sep=1pt,align=right},
    val/.style={anchor=west,inner sep=2pt,font=\footnotesize\bfseries},
  }
  % Solid bar from the 1x baseline to lo; a lighter extension covers lo..hi
  % when a row summarizes several cells. Values sit right of the bar end, or
  % right of the baseline when the bar points left.
  \newcommand{\ratiobar}[6]{% #1 y, #2 lo, #3 hi, #4 style, #5 label, #6 value text
    \pgfmathsetmacro{\xa}{\s*log2(#2/0.5)}
    \pgfmathsetmacro{\xb}{\s*log2(#3/0.5)}
    \node[lbl] at (-0.15,#1) {#5};
    \ifdim #3pt>#2pt
      \fill[#4,opacity=0.35] (\xa,#1-\barh/2) rectangle (\xb,#1+\barh/2);
    \fi
    \ifdim #2pt<1pt
      \fill[#4] (\xa,#1-\barh/2) rectangle (\xone,#1+\barh/2);
    \else
      \fill[#4] (\xone,#1-\barh/2) rectangle (\xa,#1+\barh/2);
    \fi
    \pgfmathsetmacro{\xv}{max(\xb,\xone)}
    \node[val] at (\xv,#1) {#6};
  }
  \newcommand{\ratioaxis}[2]{% #1 top y, #2 bottom y
    \pgfmathsetmacro{\ga}{\s*log2(0.95/0.5)}
    \pgfmathsetmacro{\gb}{\s*log2(1.05/0.5)}
    \fill[black!7] (\ga,#2) rectangle (\gb,#1);
    \foreach \t/\tl in {0.5/0.5,1/1,2/2,4/4,8/8,16/16,32/32} {
      \pgfmathsetmacro{\tx}{\s*log2(\t/0.5)}
      \draw[black!15] (\tx,#2) -- (\tx,#1);
      \node[anchor=north,text=black!55,font=\scriptsize] at (\tx,#2) {\tl$\times$};
    }
    \draw[black!45] (\xone,#2) -- (\xone,#1);
  }

  % ---- (a) RustFS: I versus M, direct S3 requests ----
  \begin{scope}
    \ratioaxis{0.45}{-7.45}
    \node[anchor=south west,font=\small\bfseries] at (-4.2,0.55)
      {(a) RustFS backend: revision I / revision M};
    \ratiobar{0}{23.9}{26.4}{gain}{Cold Range, $B=64\KiB$}{23.9--26.4$\times$}
    \ratiobar{-1}{4.15}{9.34}{gain}{Cold Range, $B=4\MiB$}{4.2--9.3$\times$}
    \ratiobar{-2}{2.08}{2.08}{gain}{Cold CPU per request$^{-1}$}{2.08$\times$}
    \ratiobar{-3}{1.504}{1.512}{gain}{Warm whole GET, 1 client}{1.50--1.51$\times$}
    \ratiobar{-4}{1.215}{1.215}{gain}{Warm CPU per request$^{-1}$}{1.22$\times$}
    \ratiobar{-5}{1.109}{1.140}{gain}{Cold whole GET}{1.11--1.14$\times$}
    \ratiobar{-6}{1.073}{1.123}{gain}{Warm whole GET, 8 clients}{1.07--1.12$\times$}
    \ratiobar{-7}{0.979}{1.512}{gain}{All 24 warm cells}{0.98--1.51$\times$}
  \end{scope}

  % ---- (b) BrewFS: range-heavy versus the default, through FUSE ----
  \begin{scope}[xshift=8.9cm]
    \ratioaxis{0.45}{-9.45}
    \node[anchor=south west,font=\small\bfseries] at (-4.2,0.55)
      {(b) \sysname{} client: range-heavy / default};
    \ratiobar{0}{23.45}{23.45}{gain}{4\,KiB random, 8 jobs, cold}{23.4$\times$}
    \ratiobar{-1}{16.8}{16.8}{gain}{64\,KiB random, 4 jobs, cold}{16.8$\times$}
    \ratiobar{-2}{3.56}{3.56}{gain}{4\,KiB random, 1 job, cold}{3.56$\times$}
    \ratiobar{-3}{1.02}{1.02}{flat}{Hot set, cold start}{1.02$\times$}
    \ratiobar{-4}{1.017}{1.017}{flat}{Mixed read/write, writes}{1.02$\times$}
    \ratiobar{-5}{1.001}{1.001}{flat}{4\,KiB random, warm}{1.00$\times$}
    \ratiobar{-6}{0.995}{0.995}{flat}{Mixed read/write, reads}{1.00$\times$}
    \ratiobar{-7}{0.98}{0.98}{flat}{1\,MiB sequential read, cold}{0.98$\times$}
    \ratiobar{-8}{0.97}{0.97}{flat}{Cold read of new file}{0.97$\times$}
    \ratiobar{-9}{0.91}{0.91}{loss}{256\,MiB sequential write}{0.91$\times$}
  \end{scope}
\end{tikzpicture}
\caption{Preliminary layout measurements on one cloud VM (EC~2+2,
four directories on one ESSD). Bars show candidate/baseline ratios on a
logarithmic axis; the gray band marks $\pm5\%$. (a)~Direct S3 requests,
eight paired cycles; where a row covers several request windows of one
kind, the lighter part spans their minimum to maximum. CPU rows show the
reciprocal of CPU time per request.
(b)~FUSE workloads; cold cells use the two cycles measured before the
environment degraded, warm and mixed cells all eight cycles.}
\label{fig:layout-preliminary}
\end{figure*}

